\documentclass[10pt,letterpaper]{article}
\usepackage[margin=0.9in]{geometry}
\usepackage{graphicx,xcolor,booktabs,enumitem,titlesec,longtable,hyperref}
\usepackage{parskip}
\hypersetup{colorlinks=true,linkcolor=blue!60!black,urlcolor=blue!60!black}
\titleformat{\section}{\large\bfseries\color{blue!50!black}}{\thesection.}{0.5em}{}
\titleformat{\subsection}{\normalsize\bfseries}{\thesubsection.}{0.5em}{}
\setlist[itemize]{leftmargin=*,topsep=2pt,itemsep=1pt}
\setlength{\parskip}{4pt}
\title{\textbf{PDE Top-30 Deep-Read Report}\\\large Stage 5: LLM-refined reproducibility analysis}
\author{OpenClaw / Ollie --- subagent for Rick}
\date{2026-04-24}
\begin{document}
\maketitle
\tableofcontents
\newpage

\section{Executive Summary}

This report distills Stage-5 deep reads of the top-30 PDE papers surviving Stage~4 signal-based scoring (all with \emph{full PDFs} in hand --- abstract-only papers excluded per Rick's instructions). For each paper, we read sections from the extracted text (methods / experiments / conclusions / tail-references), verified advertised GitHub/GitLab URLs by HTTP check, and produced a structured reproduction plan.

\subsection*{Tier list}
\begin{itemize}
\item \textbf{Priority A (start today, 11 papers):} Public, live code + clear benchmarks. Low friction, high confidence of quantitative replication.
\item \textbf{Priority B (plan next 1--2 weeks, 10 papers):} Partial code, heavy dependencies, or tractable reimplementation.
\item \textbf{Priority C (aspirational / deferred, 9 papers):} No public code or theory-focused; reimplementation cost dominates.
\end{itemize}

\subsection*{Top-10 revised ranking (by reproduction feasibility $\times$ scientific impact)}
\begin{enumerate}
\item \textbf{FEM vs PINNs} (Grossmann 2023) --- A, feasibility 10/10, 4h
\item \textbf{ML-accelerated CFD / jax-cfd} (Kochkov 2021) --- A, 9/10, 6h, 1123 cites
\item \textbf{Dedalus} (Burns 2019) --- A, 9/10, 4h, 485 cites
\item \textbf{Latent Spectral Models} (Wu 2023) --- A, 9/10, 6h
\item \textbf{Walk on Stars} (Sawhney 2023) --- A, 9/10, 5h
\item \textbf{Koopman Neural Operator} (Xiong 2023) --- A, 8/10, 6h
\item \textbf{Laplace Neural Operator} (Cao 2023) --- A, 8/10, 4h
\item \textbf{Fast Poisson spectral} (Fortunato/Townsend 2017) --- A, 8/10, 4h
\item \textbf{Lightning Laplace/Helmholtz} (Gopal/Trefethen 2019) --- A, 9/10, 3h
\item \textbf{Kinetic.jl} (Xiao 2021) --- A, 9/10, 4h
\end{enumerate}

\subsection*{Key changes from Stage-4 rank}
\begin{itemize}
\item \textbf{Promoted to A:} Kochkov (up from 6), Dedalus (up from 7) --- both have exemplary public code.
\item \textbf{Demoted to C:} Kopanicakova 2023 (code ``upon acceptance'' only), Foucart 2022 (github.com/corbett5/amr\_dl returns 404), Bessemoulin 2022 / Blömker 2013 / Gantner 2020 (all theory-only, no code).
\item \textbf{Off-topic flag:} Xu 2022 ``Poisson Flow Generative Models'' is a generative ML paper that uses the Poisson equation as an \emph{analogy}; it's not really a PDE-solver paper. Retained as B-tier since code is excellent but scope is a mismatch.
\end{itemize}

\subsection*{Domain / method distribution}
\begin{itemize}
\item ML-operator / PINN: 10 papers (\textasciitilde 33\%)
\item Classical spectral / FD / FV: 8 papers
\item Preconditioning / DD / multigrid: 4 papers
\item Stochastic / kinetic / fractional: 5 papers
\item Monte Carlo / grid-free: 2 papers
\item AMR / adaptivity: 3 papers
\end{itemize}

A healthy 40--60 split between ML-forward and classical numerics, matching our strategic intent.

\subsection*{Recommended ``start today'' set}
Spawn Wave~1 immediately (5 parallel subagents, $\sim$8h wallclock): FEM-vs-PINNs, Walk on Stars, Fast Poisson Spectral, Lightning Laplace, Kinetic.jl. See \texttt{top30\_batch\_schedule.md} for full 15-paper waves.

\subsection*{Cross-check against existing replications}
A text search across \texttt{\textasciitilde/Dropbox/REPLICATE-PROJECT/} found \emph{no} matching titles or DOIs among the existing $\sim 35$ subdirectories (which cover cosmology, materials, biology, NLP, etc.). All 30 PDE candidates are novel to our replication portfolio.

\newpage
\section{Per-paper deep reads}

\subsection*{\#1. Can physics-informed neural networks beat the finite element method?}
\textbf{Authors}: T. G. Grossmann;U. J. Komorowska;J. Latz;C. Schönlieb \textbullet{} \textbf{Year}: 2023 \textbullet{} \textbf{Cites}: 188 \textbullet{} \textbf{DOI}: \texttt{10.1093/imamat/hxae011} \\
\textbf{Priority}: A \textbullet{} \textbf{Revised score}: 19/20 \textbullet{} \textbf{Topic}: ML-operator/PINN

\textbf{Problem.} Systematically compare PINNs versus the finite element method (FEM) for solving a suite of linear and nonlinear PDEs. The question: can PINNs actually beat FEM in accuracy and runtime?

\textbf{Approach.} Solve 1D/2D/3D Poisson, 1D Allen-Cahn, and 1D/2D semilinear Schrödinger with both standard Galerkin FEM (FEniCS) and PINNs (PyTorch). Sweep PINN architectures (layers, nodes) and FEM mesh sizes. Measure relative error on independent test grid, solution-time and evaluation-time.

\textbf{Key results.}
\begin{itemize}
\item For all PDEs tested, FEM achieved same-or-better accuracy at 2–3 orders of magnitude lower cost than PINNs.
\item PINNs were only competitive on pointwise evaluation cost after training, not on total time to solution.
\item Larger PINN networks did not yield monotonic accuracy gains; width mattered more than depth.
\end{itemize}

\textbf{Main contribution.} Carefully-controlled empirical disproof of the common claim that PINNs outperform FEM for standard elliptic/parabolic PDEs.

\textbf{Reproducibility.} Code available: \emph{yes}. URL: \texttt{https://github.com/TamaraGrossmann/FEM-vs-PINNs (HTTP 200)}. Language: python (FEniCS + PyTorch). Completeness: 9/10. Data: Analytical/manufactured solutions, generated in-repo.

\textbf{Benchmarks.} Poisson 1D/2D/3D, Allen-Cahn 1D, Schrödinger 1D/2D (semilinear)

\textbf{Reproduction plan} (feasibility 10/10, \textasciitilde4h, compute: 1 GPU):
\begin{itemize}
\item 1. Clone github.com/TamaraGrossmann/FEM-vs-PINNs and install FEniCS + PyTorch.
\item 2. Run 1D Poisson FEM (their fem\_poisson1d.py) across mesh refinements; record L2 relative errors vs DOF.
\item 3. Train PINNs at specified architectures on 1D Poisson; record train time and L2 error.
\item 4. Reproduce Pareto plot (Fig 2): FEM dominates PINN in time-accuracy tradeoff.
\item 5. Repeat for Allen-Cahn 1D and confirm qualitative conclusion.
\item 6. Optional: extend with 3D Poisson to reproduce the one regime (eval time) where PINN wins.
\end{itemize}
\textbf{Expected fidelity:} Quantitative match within \~{}10\% on relative-error tables; qualitative reproduction of all Pareto figures.. \textbf{Known risks:} PINN training is stochastic; averaged over several seeds for stable numbers.; FEniCS version drift (legacy vs FEniCSx)..

\textbf{Follow-on questions.}
\begin{itemize}
\item Does the PINN-vs-FEM gap close at higher spatial dimensions (d≥5)?
\item How do adaptive activation functions or Fourier feature embeddings change the Pareto?
\item Is the conclusion robust under parametric PDE families where FEM must be resolved many times (surrogate regime)?
\item What happens under irregular geometry where FEM meshing is itself expensive?
\item Does causal PINN training alter the Allen-Cahn verdict?
\end{itemize}

\textbf{Rationale for priority:} \emph{Perfect match to our needs: open code, FEniCS+PyTorch only, direct quantitative comparison. Keep top rank.}

\medskip\hrule\medskip

\subsection*{\#2. Solving High-Dimensional PDEs with Latent Spectral Models}
\textbf{Authors}: Haixu Wu;Tengge Hu;Huakun Luo;Jianmin Wang;Mingsheng Long \textbullet{} \textbf{Year}: 2023 \textbullet{} \textbf{Cites}: 95 \textbullet{} \textbf{DOI}: \texttt{10.48550/arXiv.2301.12664} \\
\textbf{Priority}: A \textbullet{} \textbf{Revised score}: 18/20 \textbullet{} \textbf{Topic}: ML-operator/PINN

\textbf{Problem.} Neural operators work in coordinate space and suffer curse-of-dimensionality for high-dim PDEs. Learn PDEs in a compact latent spectral space.

\textbf{Approach.} Latent Spectral Models (LSM) use hierarchical projection to a low-dim latent, then a 'neural spectral block' that approximates the operator with trainable basis functions. Theoretical guarantees via spectral approximation theory.

\textbf{Key results.}
\begin{itemize}
\item SOTA on 7 solid+fluid benchmarks (Elasticity, Plasticity, Navier-Stokes, Darcy, Airfoil, Pipe).
\item Stable performance across resolutions 32x32 → 1024x1024 on Darcy.
\item Consistent outperformance of FNO, U-FNO, WMT, Galerkin-Transformer baselines.
\end{itemize}

\textbf{Main contribution.} Operator learning in a learned latent spectral basis—decoupling representation from output grid.

\textbf{Reproducibility.} Code available: \emph{yes}. URL: \texttt{https://github.com/thuml/Latent-Spectral-Models (HTTP 200; also in thuml/Neural-Solver-Library)}. Language: python (PyTorch). Completeness: 10/10. Data: PDEBench / Takamoto benchmarks; dataset download links in repo README.

\textbf{Benchmarks.} Elasticity-P, Elasticity-G, Plasticity, Navier-Stokes (2D), Darcy, Airfoil, Pipe

\textbf{Reproduction plan} (feasibility 9/10, \textasciitilde6h, compute: 1 GPU (A100 ideal, any modern GPU sufficient)):
\begin{itemize}
\item 1. Clone thuml/Latent-Spectral-Models; install PyTorch per requirements.txt.
\item 2. Download Darcy 2D dataset (small benchmark, fast iteration).
\item 3. Train LSM-S on Darcy (should take \~{}1–2h on a single GPU); reproduce relative L2 near 0.0065 (paper value).
\item 4. Train FNO baseline from the same repo on same data; compare to LSM.
\item 5. Run resolution test across 64/128/256 on Darcy to replicate resolution-stability figure.
\item 6. (Optional) Navier-Stokes 2D benchmark for a harder check.
\end{itemize}
\textbf{Expected fidelity:} Quantitative match within 5–10\% on Darcy \& NS-2D L2 errors.. \textbf{Known risks:} GPU memory at 1024x1024 resolution.; Dataset download bandwidth (10+ GB)..

\textbf{Follow-on questions.}
\begin{itemize}
\item How does LSM perform on 3D turbulence datasets (Johns Hopkins Turbulence Database)?
\item Can the latent-spectral basis be made physics-aware (e.g., divergence-free)?
\item What is the sample-efficiency curve vs. FNO?
\item Does LSM generalize zero-shot across parametric PDE families?
\item Can the approach be combined with multigrid for hybrid solver acceleration?
\end{itemize}

\textbf{Rationale for priority:} \emph{Well-supported repo and benchmark datasets; strong baselines. Prime candidate.}

\medskip\hrule\medskip

\subsection*{\#3. Koopman neural operator as a mesh-free solver of non-linear partial differential equations}
\textbf{Authors}: Wei Xiong;Xiaomeng Huang;Ziyang Zhang;Ruixuan Deng;Pei Sun;Yang Tian \textbullet{} \textbf{Year}: 2023 \textbullet{} \textbf{Cites}: 74 \textbullet{} \textbf{DOI}: \texttt{10.48550/arXiv.2301.10022} \\
\textbf{Priority}: A \textbullet{} \textbf{Revised score}: 17/20 \textbullet{} \textbf{Topic}: ML-operator/PINN

\textbf{Problem.} Existing neural operators become inaccurate and less interpretable for long-term/non-linear PDE evolution.

\textbf{Approach.} Koopman Neural Operator (KNO) formulates the PDE solution as a non-linear dynamical system and approximates its Koopman operator (infinite-dim linear) acting on flow embeddings. Uses an FNO-like lifting but with a Koopman block.

\textbf{Key results.}
\begin{itemize}
\item Outperforms FNO on mesh-independent Burgers experiments and long-term prediction on 5 PDEs incl. 2D Navier-Stokes and Rayleigh-Bénard.
\item Zero-shot prediction across untrained discretization granularities.
\item Applied to real systems: global water vapor, western boundary currents.
\end{itemize}

\textbf{Main contribution.} First Koopman-theoretic neural operator with demonstrated long-horizon stability on turbulent PDEs.

\textbf{Reproducibility.} Code available: \emph{yes}. URL: \texttt{https://github.com/Koopman-Laboratory/KoopmanLab (HTTP 200)}. Language: python (PyTorch). Completeness: 9/10. Data: Benchmark datasets reused from FNO paper (Burgers, NS-2D) + Copernicus ocean data.

\textbf{Benchmarks.} Bateman-Burgers 1D, 2D Navier-Stokes, Rayleigh-Bénard, Shallow Water, Diffusion-reaction

\textbf{Reproduction plan} (feasibility 8/10, \textasciitilde6h, compute: 1 GPU (A100 40GB used by authors)):
\begin{itemize}
\item 1. Clone Koopman-Laboratory/KoopmanLab; install via pip install -e .
\item 2. Download FNO's Burgers dataset; run python train\_burgers.py with KNO config.
\item 3. Verify reported relative L2 \~{} reported value (check paper Table 2).
\item 4. Train FNO baseline for head-to-head.
\item 5. Extend to NS-2D long-term (100+ steps) and check stability.
\item 6. Zero-shot: train at res=32, evaluate at res=64/128.
\end{itemize}
\textbf{Expected fidelity:} Within 10-15\% on L2; qualitative reproduction of Koopman-eigenvalue analysis.. \textbf{Known risks:} Long-term rollout stability is finicky.; Possible version drift in KoopmanLab package..

\textbf{Follow-on questions.}
\begin{itemize}
\item How does KNO compare against POD-DL-ROM on same benchmarks?
\item Is the learned Koopman spectrum interpretable / linked to physical modes?
\item Can KNO handle boundary changes (mesh-free in a stronger sense)?
\item Performance on 3D turbulence?
\item Hybrid KNO + classical time-integrator for stability guarantees?
\end{itemize}

\textbf{Rationale for priority:} \emph{Strong code, clear benchmarks. Slightly trickier than LSM due to long-rollout stability.}

\medskip\hrule\medskip

\subsection*{\#4. Lifex-cfd: an Open-source Computational Fluid Dynamics Solver for Cardiovascular Applications}
\textbf{Authors}: P. C. Africa;Ivan Fumagalli;M. Bucelli;Alberto Zingaro;L. Dede’;A. Quarteroni \textbullet{} \textbf{Year}: 2023 \textbullet{} \textbf{Cites}: 61 \textbullet{} \textbf{DOI}: \texttt{10.48550/arXiv.2304.12032} \\
\textbf{Priority}: B \textbullet{} \textbf{Revised score}: 16/20 \textbullet{} \textbf{Topic}: Fluid/NavierStokes

\textbf{Problem.} Open-source CFD solver specialized for cardiovascular simulations (blood flow in moving cardiac chambers + valves).

\textbf{Approach.} Finite-element incompressible Navier-Stokes solver on deal.II/PETSc. ALE for moving domains, VMS-LES for transitional turbulence, RIIS for immersed valves. SUPG-PSPG stabilization.

\textbf{Key results.}
\begin{itemize}
\item Verified ideal convergence orders on Beltrami benchmark.
\item Near-ideal parallel speedup on Taylor-Green vortex.
\item Patient-specific cases for idealized and real vasculature.
\end{itemize}

\textbf{Main contribution.} Production-quality, LGPL-v3 cardiovascular CFD with careful verification and parallel scalability.

\textbf{Reproducibility.} Code available: \emph{yes}. URL: \texttt{https://gitlab.com/lifex/lifex-cfd (HTTP 200) + Zenodo 10.5281/zenodo.7852088 for input data}. Language: C++ (deal.II, PETSc, Trilinos). Completeness: 10/10. Data: Zenodo 7852088 - parameter files, meshes, auxiliary data.

\textbf{Benchmarks.} Beltrami flow, Taylor-Green vortex, patient-specific aorta/LV

\textbf{Reproduction plan} (feasibility 7/10, \textasciitilde10h, compute: multi-CPU node; HPC for full cardiac test):
\begin{itemize}
\item 1. Build deal.II with MPI/PETSc/Trilinos; build lifex core, then lifex-cfd.
\item 2. Download Zenodo 7852088; run Beltrami benchmark → verify convergence table.
\item 3. Run Taylor-Green vortex on 32-core node; verify expected kinetic-energy dissipation curve.
\item 4. (Optional) Run one patient-specific idealized aorta case from the provided data.
\end{itemize}
\textbf{Expected fidelity:} Quantitative match on convergence rates and kinetic-energy dissipation curves.. \textbf{Known risks:} deal.II build is non-trivial (1–3h compile).; Full cardiac cases may require HPC queue time..

\textbf{Follow-on questions.}
\begin{itemize}
\item Can the lifex-cfd Navier-Stokes solver be coupled with Windkessel BCs from the literature?
\item Performance vs. OpenFOAM on the same Beltrami benchmark?
\item Can ML turbulence models replace VMS-LES?
\item GPU port via deal.II's matrix-free infrastructure?
\item Verification on pulsatile non-Newtonian rheology.
\end{itemize}

\textbf{Rationale for priority:} \emph{Excellent reproducibility but heavy dependency stack (deal.II + Trilinos) inflates setup time. Great if we want a cardiovascular reference; not a 'start today' pick.}

\medskip\hrule\medskip

\subsection*{\#5. Enhancing training of physics-informed neural networks using domain-decomposition based preconditioning strategies}
\textbf{Authors}: Alena Kopanicáková;Hardik Kothari;G. Karniadakis;Rolf Krause \textbullet{} \textbf{Year}: 2023 \textbullet{} \textbf{Cites}: 28 \textbullet{} \textbf{DOI}: \texttt{10.1137/23M1583375} \\
\textbf{Priority}: C \textbullet{} \textbf{Revised score}: 13/20 \textbullet{} \textbf{Topic}: ML-operator/PINN

\textbf{Problem.} PINN training often converges slowly with L-BFGS; preconditioning is underexplored.

\textbf{Approach.} Introduce Schwarz-domain-decomposition-based nonlinear preconditioners (additive SPQN and multiplicative MPQN) that decompose network parameters layer-wise and apply L-BFGS within blocks.

\textbf{Key results.}
\begin{itemize}
\item Significant convergence acceleration vs vanilla L-BFGS/Adam on several PINN benchmarks.
\item Additive preconditioner is embarrassingly parallel → new model-parallelism route for PINNs.
\item Lower achieved relative error at fixed training budget.
\end{itemize}

\textbf{Main contribution.} Brings classical domain-decomposition preconditioning into nonlinear optimization for PINNs.

\textbf{Reproducibility.} Code available: \emph{no}. URL: \texttt{Promised upon acceptance (Bitbucket/alena\_ referenced in text); SIAM J. Sci. Comput. 2023 version should have code appendix}. Language: python (PyTorch per conventions of Karniadakis group). Completeness: 4/10. Data: Manufactured solutions for Burgers, Allen-Cahn, advection tests.

\textbf{Benchmarks.} Burgers, Allen-Cahn, advection-diffusion

\textbf{Reproduction plan} (feasibility 6/10, \textasciitilde8h, compute: 1 GPU):
\begin{itemize}
\item 1. Implement layer-wise Schwarz block L-BFGS from scratch guided by Algorithm 1 in paper.
\item 2. Reproduce Burgers PINN baseline with Adam and plain L-BFGS.
\item 3. Add additive SPQN wrapper; verify convergence curves reported in Fig 5.
\item 4. Measure final Erel and compare to paper Tables 3-5.
\end{itemize}
\textbf{Expected fidelity:} Qualitative reproduction of convergence acceleration; numerical match within 20\% given stochasticity.. \textbf{Known risks:} No public code yet → full re-implementation required.; L-BFGS line search finicky in Schwarz setting..

\textbf{Follow-on questions.}
\begin{itemize}
\item Does the preconditioner help for stiff PINN losses (low-viscosity Burgers)?
\item How does it compose with causal/self-adaptive PINN weighting?
\item Scaling to 100-layer PINNs?
\item Comparison to natural-gradient PINN variants?
\item Theoretical convergence rate of the nonlinear Schwarz iteration.
\end{itemize}

\textbf{Rationale for priority:} \emph{Strong idea, but code not public at time of writing — demotes from A to C until repo is released or full reimplementation budget allotted.}

\medskip\hrule\medskip

\subsection*{\#6. Machine learning–accelerated computational fluid dynamics}
\textbf{Authors}: Dmitrii Kochkov;Jamie A. Smith;Ayya Alieva;Qing Wang;M. Brenner;Stephan Hoyer \textbullet{} \textbf{Year}: 2021 \textbullet{} \textbf{Cites}: 1123 \textbullet{} \textbf{DOI}: \texttt{10.1073/pnas.2101784118} \\
\textbf{Priority}: A \textbullet{} \textbf{Revised score}: 19/20 \textbullet{} \textbf{Topic}: ML-operator/PINN

\textbf{Problem.} DNS of Navier-Stokes is limited by smallest-scale resolution; LES relies on hand-tuned closures.

\textbf{Approach.} End-to-end differentiable CFD where a CNN corrects the convective flux on a coarse grid; trained against high-res DNS of 2D Kolmogorov flow with model-consistent rollouts (32 steps).

\textbf{Key results.}
\begin{itemize}
\item Matches DNS accuracy at 8–10× coarser grid → \~{}40–80× compute speedup.
\item Stable long-horizon rollouts (tens of thousands of time steps).
\item Outperforms LES closures at equivalent cost.
\end{itemize}

\textbf{Main contribution.} First large-scale demonstration of ML-accelerated CFD with DNS-fidelity on 2D turbulence.

\textbf{Reproducibility.} Code available: \emph{yes}. URL: \texttt{https://github.com/google/jax-cfd (HTTP 200)}. Language: python (JAX). Completeness: 10/10. Data: Datasets generated in-repo (seeded); also published via TFDS.

\textbf{Benchmarks.} 2D Kolmogorov flow Re=1000, 2D decaying turbulence

\textbf{Reproduction plan} (feasibility 9/10, \textasciitilde6h, compute: 1 GPU (Spark or A100)):
\begin{itemize}
\item 1. pip install jax-cfd; run the notebook `notebooks/kolmogorov\_flow.ipynb`.
\item 2. Generate ground-truth DNS at 2048² (may need to reduce to 512² for time budget).
\item 3. Train the learned correction model (default config).
\item 4. Compare vorticity correlation vs baseline DNS at same coarse resolution.
\item 5. Reproduce Fig 1 quantitative results: coarse-DNS vs ML-DNS vorticity correlations.
\end{itemize}
\textbf{Expected fidelity:} Quantitative match on vorticity correlation and energy spectra.. \textbf{Known risks:} DNS at 2048² is expensive; may need to cap at 1024².; JAX-cfd's default training can take hours..

\textbf{Follow-on questions.}
\begin{itemize}
\item Does the same approach work for 3D turbulence?
\item Transfer learning across Reynolds numbers?
\item Hybrid with entropy-stable limiters for shocks?
\item Training with less high-res data (sample efficiency)?
\item Integrating this with weather model cores (NWP)?
\end{itemize}

\textbf{Rationale for priority:} \emph{1123 citations, Google JAX-CFD repo is exemplary, data generation built in. Upgrade to co-\#1 priority.}

\medskip\hrule\medskip

\subsection*{\#7. Dedalus: A flexible framework for numerical simulations with spectral methods}
\textbf{Authors}: K. Burns;G. Vasil;J. Oishi;D. Lecoanet;B. Brown \textbullet{} \textbf{Year}: 2019 \textbullet{} \textbf{Cites}: 485 \textbullet{} \textbf{DOI}: \texttt{10.1103/PhysRevResearch.2.023068} \\
\textbf{Priority}: A \textbullet{} \textbf{Revised score}: 18/20 \textbullet{} \textbf{Topic}: Fluid/NavierStokes

\textbf{Problem.} Need a flexible framework to solve general PDEs spectrally in many geometries (Cartesian, spherical, annular).

\textbf{Approach.} Dedalus: object-oriented Python framework representing PDEs symbolically; uses modified sparse tau method with Chebyshev/Fourier bases; banded matrices via Dirichlet preconditioning. Parallel via MPI.

\textbf{Key results.}
\begin{itemize}
\item Demonstrated on Kelvin-Helmholtz, Rayleigh-Bénard, Kolmogorov flow, self-gravitating cylinders.
\item Banded systems give O(N) solves per timestep.
\item Scalable to O(10k) cores on ALCF/NERSC.
\end{itemize}

\textbf{Main contribution.} Widely-adopted (485 cites) production-grade spectral PDE framework.

\textbf{Reproducibility.} Code available: \emph{yes}. URL: \texttt{https://github.com/DedalusProject/dedalus (HTTP 200) + DedalusProject/methods\_paper\_examples}. Language: python (NumPy, MPI, h5py). Completeness: 10/10. Data: Example scripts in methods\_paper\_examples repo.

\textbf{Benchmarks.} Poisson 1D/2D, Helmholtz, Kelvin-Helmholtz 2D, Rayleigh-Bénard, reaction-diffusion

\textbf{Reproduction plan} (feasibility 9/10, \textasciitilde4h, compute: CPU (multi-core); MPI for 2D+ cases):
\begin{itemize}
\item 1. conda install -c conda-forge dedalus.
\item 2. Clone DedalusProject/methods\_paper\_examples.
\item 3. Run 1D Poisson example; verify spectral convergence.
\item 4. Run 2D Kelvin-Helmholtz example; compare to published vorticity field.
\item 5. Run Rayleigh-Bénard at Ra=1e6; check Nusselt number scaling.
\end{itemize}
\textbf{Expected fidelity:} Exact match on spectral convergence; qualitative on turbulent cases.. \textbf{Known risks:} Dedalus v3 changed API significantly; methods\_paper\_examples may target v2.; MPI build can be finicky..

\textbf{Follow-on questions.}
\begin{itemize}
\item Benchmark Dedalus vs Julia's ApproxFun + FourierFlows on same problems.
\item GPU port via JAX or CUDA?
\item Automatic-differentiation overlay for inverse problems?
\item Use for dynamo simulations — replicate Vasil's ICC benchmark.
\item Coupling with AMR for local resolution.
\end{itemize}

\textbf{Rationale for priority:} \emph{Mature, heavily cited, easy install. Upgrade — excellent starting point.}

\medskip\hrule\medskip

\subsection*{\#8. Physics-informed neural networks for solving Reynolds-averaged Navier-Stokes equations}
\textbf{Authors}: Hamidreza Eivazi;Mojtaba Tahani;P. Schlatter;R. Vinuesa \textbullet{} \textbf{Year}: 2021 \textbullet{} \textbf{Cites}: 404 \textbullet{} \textbf{DOI}: \texttt{10.1063/5.0095270} \\
\textbf{Priority}: B \textbullet{} \textbf{Revised score}: 13/20 \textbullet{} \textbf{Topic}: Fluid/NavierStokes

\textbf{Problem.} Turbulence closure is hard; can PINNs learn RANS solutions using only boundary data?

\textbf{Approach.} Fully-connected PINN taking (x,y) → (U,V,P,u²,v²,uv). Loss = PDE residual of RANS + data match at boundaries (no interior closure assumption).

\textbf{Key results.}
\begin{itemize}
\item <1\% error on laminar Falkner-Skan.
\item Good accuracy on ZPG/APG boundary layers and NACA-0012 airfoil without a turbulence model.
\item Reynolds stresses recovered consistently.
\end{itemize}

\textbf{Main contribution.} Shows PINNs can bypass explicit turbulence closure if boundary Reynolds-stress data are provided.

\textbf{Reproducibility.} Code available: \emph{no}. URL: \texttt{None advertised; may be available upon request from authors}. Language: python (TensorFlow per text). Completeness: 3/10. Data: DNS/LES reference data for ZPG/APG BL and NACA (public DNS databases).

\textbf{Benchmarks.} Falkner-Skan, ZPG BL, APG BL, NACA-0012

\textbf{Reproduction plan} (feasibility 6/10, \textasciitilde8h, compute: 1 GPU):
\begin{itemize}
\item 1. Set up a 5-output PINN in DeepXDE with RANS residual.
\item 2. Generate Falkner-Skan analytical solution; train; expect <1\% error.
\item 3. Use Jiménez DNS for ZPG BL as boundary data; train.
\item 4. Compare predicted U,V,P,u²,v²,uv against DNS at interior stations.
\end{itemize}
\textbf{Expected fidelity:} Qualitative (shapes of U,V profiles); quantitative L2 within \~{}10\% of paper.. \textbf{Known risks:} No public code → reimplementation.; PINN for turbulent cases is notoriously sensitive to loss-weight scheduling..

\textbf{Follow-on questions.}
\begin{itemize}
\item Does it extend to 3D separated flows?
\item Comparison against data-driven SA or k-omega closures?
\item Uncertainty quantification on the recovered Reynolds stresses.
\item Integration with real PIV data instead of DNS.
\item Does causal/time-marching PINN improve on this steady formulation?
\end{itemize}

\textbf{Rationale for priority:} \emph{Clear, impactful paper but no public code lowers feasibility; still a B-tier candidate given DeepXDE reimplementation is tractable.}

\medskip\hrule\medskip

\subsection*{\#9. Learning in Modal Space: Solving Time-Dependent Stochastic PDEs Using Physics-Informed Neural Networks}
\textbf{Authors}: Dongkun Zhang;Ling Guo;G. Karniadakis \textbullet{} \textbf{Year}: 2019 \textbullet{} \textbf{Cites}: 274 \textbullet{} \textbf{DOI}: \texttt{10.1137/19m1260141} \\
\textbf{Priority}: C \textbullet{} \textbf{Revised score}: 11/20 \textbullet{} \textbf{Topic}: ML-operator/PINN

\textbf{Problem.} Long-time integration of stochastic PDEs with arbitrary initial data — classical DO/BO methods have singular covariance and eigenvalue-crossing issues.

\textbf{Approach.} NN-DO and NN-BO: parameterize modes and coefficients by neural networks; enforce DO/BO constraints implicitly in the loss (no explicit time derivatives of modes).

\textbf{Key results.}
\begin{itemize}
\item Works when original DO fails (deterministic IC case).
\item RMSE for BO random coefficients \~{}0.05–0.29 at early/late times.
\item Handles long-time integration via subdomain-in-time decomposition.
\end{itemize}

\textbf{Main contribution.} Trainable DO/BO decomposition avoiding the pathological assumptions of classical versions.

\textbf{Reproducibility.} Code available: \emph{no}. URL: \texttt{None advertised; possibly available from Karniadakis group on request}. Language: python (TensorFlow 1.x, typical of 2019 Karniadakis). Completeness: 3/10. Data: Manufactured SPDE problems (linear stochastic advection, stochastic diffusion-reaction, etc.).

\textbf{Benchmarks.} Stochastic advection 1D, Stochastic diffusion-reaction 1D, Stochastic Burgers

\textbf{Reproduction plan} (feasibility 5/10, \textasciitilde10h, compute: 1 GPU):
\begin{itemize}
\item 1. Reimplement the DO/BO loss with NN parameterization (algorithms 1-2 in paper).
\item 2. Benchmark on stochastic advection; match RMSE \~{}0.05 after 1000 iterations.
\item 3. Extend to diffusion-reaction; match Table 5 random-coefficient RMSE.
\end{itemize}
\textbf{Expected fidelity:} Qualitative match; quantitative within 2× of reported RMSE.. \textbf{Known risks:} No code; 2019 paper → reimplementation.; Subdomain time-splitting detail is fragile..

\textbf{Follow-on questions.}
\begin{itemize}
\item Extension to stochastic Navier-Stokes 2D?
\item Adaptive basis dimension during training?
\item Comparison to stochastic Galerkin + PC expansion.
\item Inverse problems under stochastic forcing.
\item How much training data does the inverse setup require?
\end{itemize}

\textbf{Rationale for priority:} \emph{Niche but no public code; demote to C.}

\medskip\hrule\medskip

\subsection*{\#10. Poisson Flow Generative Models}
\textbf{Authors}: Yilun Xu;Ziming Liu;M. Tegmark;T. Jaakkola \textbullet{} \textbf{Year}: 2022 \textbullet{} \textbf{Cites}: 115 \textbullet{} \textbf{DOI}: \texttt{10.48550/arXiv.2209.11178} \\
\textbf{Priority}: B \textbullet{} \textbf{Revised score}: 14/20 \textbullet{} \textbf{Topic}: Elliptic

\textbf{Problem.} Normalizing flows / SDE-based generative models have slow sampling and mode-collapse issues.

\textbf{Approach.} Interpret data as charges in z=0 plane of a (d+1)-dim space; the gradient of the Poisson equation solution induces a deterministic flow that maps data → uniform hemisphere. Backward ODE gives sampling.

\textbf{Key results.}
\begin{itemize}
\item Inception 9.68, FID 2.35 on CIFAR-10 (SOTA at time among normalizing flows).
\item 10-20× faster sampling than VP/VE-SDE baselines at equal FID.
\item Robust to network choice and step size.
\end{itemize}

\textbf{Main contribution.} A PDE-physics-inspired generative model with fewer sampling steps and good FID.

\textbf{Reproducibility.} Code available: \emph{yes}. URL: \texttt{https://github.com/Newbeeer/poisson\_flow (HTTP 200)}. Language: python (PyTorch). Completeness: 10/10. Data: CIFAR-10, CelebA (public).

\textbf{Benchmarks.} CIFAR-10 FID/IS, CelebA 64x64, LSUN bedroom

\textbf{Reproduction plan} (feasibility 9/10, \textasciitilde10h, compute: 1 GPU (A100 ideal); training 1-2 days):
\begin{itemize}
\item 1. Clone Newbeeer/poisson\_flow.
\item 2. Use provided CIFAR-10 config; train.
\item 3. Evaluate FID/IS using their eval script.
\item 4. Reproduce sampling-speedup table (Table in paper).
\end{itemize}
\textbf{Expected fidelity:} FID within \~{}0.2 of published 2.35.. \textbf{Known risks:} Training takes \~{}1 day.; FID calculation depends on reference stats; ensure matching CIFAR-10 stats file..

\textbf{Follow-on questions.}
\begin{itemize}
\item How does PFGM compare to later PFGM++ and diffusion-transformer models?
\item Does the Poisson analogy extend to discrete data (text)?
\item Connection to score matching and Wasserstein flows.
\item Ablate the hemisphere geometry — is it necessary?
\item Large-scale text-to-image applications.
\end{itemize}

\textbf{Rationale for priority:} \emph{Not really a PDE-solver paper (it's a generative model using Poisson as analogy). Good code but off-topic for PDE replication; B tier at best for 'PDE' bucket.}

\medskip\hrule\medskip

\subsection*{\#11. LNO: Laplace Neural Operator for Solving Differential Equations}
\textbf{Authors}: Qianying Cao;S. Goswami;G. Karniadakis \textbullet{} \textbf{Year}: 2023 \textbullet{} \textbf{Cites}: 56 \textbullet{} \textbf{DOI}: \texttt{10.48550/arXiv.2303.10528} \\
\textbf{Priority}: A \textbullet{} \textbf{Revised score}: 16/20 \textbullet{} \textbf{Topic}: ML-operator/PINN

\textbf{Problem.} FNO handles only periodic signals and cannot represent non-periodic transient responses well.

\textbf{Approach.} Laplace Neural Operator (LNO): replace Fourier transform block with Laplace transform (pole-residue representation). Natural handling of transients via poles.

\textbf{Key results.}
\begin{itemize}
\item LNO with 1 layer beats FNO with 4 layers on 3 ODEs (Duffing, driven pendulum, Lorenz) and 3 PDEs (Euler-Bernoulli beam, diffusion, reaction-diffusion).
\item Exponential convergence claimed for linear problems.
\item Better transient-response capture in undamped cases.
\end{itemize}

\textbf{Main contribution.} Laplace-transform-based neural operator handling non-periodic and transient signals.

\textbf{Reproducibility.} Code available: \emph{yes}. URL: \texttt{https://github.com/qianyingcao/Laplace-Neural-Operator (HTTP 200)}. Language: python (PyTorch). Completeness: 9/10. Data: Generated in-repo via SciPy ODE solvers and analytical solutions.

\textbf{Benchmarks.} Duffing, driven pendulum, Lorenz, Euler-Bernoulli beam, diffusion, reaction-diffusion

\textbf{Reproduction plan} (feasibility 8/10, \textasciitilde4h, compute: 1 GPU):
\begin{itemize}
\item 1. Clone qianyingcao/Laplace-Neural-Operator.
\item 2. Run Duffing training; verify relative L2 < FNO baseline.
\item 3. Run reaction-diffusion; verify Table 2 values.
\item 4. Sweep \# poles; produce the convergence-vs-poles plot.
\end{itemize}
\textbf{Expected fidelity:} Quantitative match on Table 2 within 10\%.. \textbf{Known risks:} Complex-valued arithmetic in pole-residue layer can have numerical issues..

\textbf{Follow-on questions.}
\begin{itemize}
\item Extension to 2D PDEs on non-periodic domains?
\item Can the learned poles be interpreted physically?
\item Compare with wavelet/DeepONet operators.
\item Efficiency vs. FNO at equal parameter count.
\item Transfer learning across physical parameters.
\end{itemize}

\textbf{Rationale for priority:} \emph{Clean scope, public code, lightweight benchmarks. Solid A pick.}

\medskip\hrule\medskip

\subsection*{\#12. Walk on Stars: A Grid-Free Monte Carlo Method for PDEs with Neumann Boundary Conditions}
\textbf{Authors}: Rohan Sawhney;Bailey Miller;Ioannis Gkioulekas;Keenan Crane \textbullet{} \textbf{Year}: 2023 \textbullet{} \textbf{Cites}: 48 \textbullet{} \textbf{DOI}: \texttt{10.1145/3592398} \\
\textbf{Priority}: A \textbullet{} \textbf{Revised score}: 17/20 \textbullet{} \textbf{Topic}: Elliptic

\textbf{Problem.} Walk-on-spheres Monte Carlo solves only Dirichlet Laplace problems; no Neumann BCs.

\textbf{Approach.} Walk on Stars (WoSt) samples on star-shaped regions whose boundaries include Neumann-boundary arcs; sampling via visibility silhouette. Solves linear elliptic PDEs with mixed BCs without meshing.

\textbf{Key results.}
\begin{itemize}
\item First grid-free MC method for Neumann conditions.
\item Works on arbitrary geometries (non-convex) and extreme geometric complexity.
\item Progressive output (like ray tracing) → useful for graphics.
\end{itemize}

\textbf{Main contribution.} Fundamentally extends WoS to the full mixed-BC regime with practical BVH-based implementation.

\textbf{Reproducibility.} Code available: \emph{yes}. URL: \texttt{https://github.com/GeometryCollective/wost-simple (HTTP 200)}. Language: C++ (with Python bindings); also simple Python reference. Completeness: 9/10. Data: Example scenes in repo.

\textbf{Benchmarks.} 2D Laplace mixed-BC, Poisson with Neumann, heat transfer (toaster-bread example)

\textbf{Reproduction plan} (feasibility 9/10, \textasciitilde5h, compute: CPU (multi-thread)):
\begin{itemize}
\item 1. Clone wost-simple; build (CMake + C++17).
\item 2. Run 2D Laplace demo with mixed Dirichlet/Neumann.
\item 3. Verify pointwise convergence rate N\^{}\{-1/2\}.
\item 4. Compare against FEniCS reference solution on same geometry.
\end{itemize}
\textbf{Expected fidelity:} Quantitative MC convergence rate match; pointwise values within 1σ of reference.. \textbf{Known risks:} Robust silhouette computation is subtle.; Limited to linear elliptic..

\textbf{Follow-on questions.}
\begin{itemize}
\item Extension to Robin BCs and anisotropic media?
\item Variance reduction via control variates with WoS.
\item Use in rendering-physics hybrid pipelines.
\item Walk on Stars for PDE-constrained optimization.
\item 3D performance on complex meshes.
\end{itemize}

\textbf{Rationale for priority:} \emph{Beautiful paper, great code. A-tier.}

\medskip\hrule\medskip

\subsection*{\#13. Fast Poisson solvers for spectral methods}
\textbf{Authors}: D. Fortunato;Alex Townsend \textbullet{} \textbf{Year}: 2017 \textbullet{} \textbf{Cites}: 47 \textbullet{} \textbf{DOI}: \texttt{10.1093/imanum/drz034} \\
\textbf{Priority}: A \textbullet{} \textbf{Revised score}: 16/20 \textbullet{} \textbf{Topic}: Elliptic

\textbf{Problem.} Fast Poisson solvers exist for FD/FE but not for spectral methods; spectral accuracy traditionally couples with O(n\^{}3) cost.

\textbf{Approach.} Ultraspherical spectral discretization → Sylvester matrix equation with 'separated spectra' → ADI method gives O(n\^{}2 log\^{}2 n) complexity. Works on square, cylinder, solid sphere, cube.

\textbf{Key results.}
\begin{itemize}
\item 100M DOF Poisson on a square solved in < 2 minutes on laptop.
\item Spectral accuracy with quasi-optimal complexity.
\item Integrated into Chebfun via chebfun2.poisson.
\end{itemize}

\textbf{Main contribution.} First optimal-complexity spectral Poisson solver via ADI on Sylvester equations.

\textbf{Reproducibility.} Code available: \emph{yes}. URL: \texttt{https://github.com/danfortunato/fast-poisson-solvers (HTTP 200) + Chebfun chebop2}. Language: MATLAB (Chebfun). Completeness: 9/10. Data: Analytical test functions in repo.

\textbf{Benchmarks.} Poisson on square, cylinder, solid sphere, cube

\textbf{Reproduction plan} (feasibility 8/10, \textasciitilde4h, compute: CPU (single laptop)):
\begin{itemize}
\item 1. Install Chebfun in MATLAB/Octave.
\item 2. u = chebfun2.poisson(f, BC, n) for n=4096 → measure time.
\item 3. Verify runtime scaling across n=64,...,16384.
\item 4. Benchmark vs. MATLAB's built-in Poisson solver on same grid.
\end{itemize}
\textbf{Expected fidelity:} Exact match on runtime scaling; residual at machine precision.. \textbf{Known risks:} MATLAB license; Octave support for Chebfun is partial..

\textbf{Follow-on questions.}
\begin{itemize}
\item Port to Python (via chebpy)?
\item Extension to variable-coefficient Poisson?
\item Biharmonic and Helmholtz analogs via same ADI approach?
\item GPU implementation?
\item Use as preconditioner in more general elliptic solvers.
\end{itemize}

\textbf{Rationale for priority:} \emph{Classic, clean, reproducible. MATLAB/Chebfun dependency is minor. A-tier.}

\medskip\hrule\medskip

\subsection*{\#14. Deep Reinforcement Learning for Adaptive Mesh Refinement}
\textbf{Authors}: C. Foucart;A. Charous;Pierre FJ Lermusiaux \textbullet{} \textbf{Year}: 2022 \textbullet{} \textbf{Cites}: 35 \textbullet{} \textbf{DOI}: \texttt{10.48550/arXiv.2209.12351} \\
\textbf{Priority}: C \textbullet{} \textbf{Revised score}: 11/20 \textbullet{} \textbf{Topic}: AMR/Multigrid

\textbf{Problem.} AMR refinement heuristics are ad-hoc; can RL learn better refinement policies?

\textbf{Approach.} Formulate AMR as POMDP with local per-element observation. Train deep Q-network to refine/coarsen HDG-FEM elements on advection and Poisson.

\textbf{Key results.}
\begin{itemize}
\item RL agent reduces DOF by \~{}20-30\% at fixed error vs heuristic.
\item Generalizes across PDEs (trained on advection, tested on Poisson).
\item Demonstrated on 1D \& 2D.
\end{itemize}

\textbf{Main contribution.} First practical RL-AMR controller shown to beat bulk-refinement heuristic on DG-FEM.

\textbf{Reproducibility.} Code available: \emph{no}. URL: \texttt{github.com/corbett5/amr\_dl (DEAD 404) — request from authors}. Language: python (custom HDG code + PyTorch/RLlib). Completeness: 3/10. Data: Manufactured solutions.

\textbf{Benchmarks.} advection 1D/2D, Poisson 1D with Gaussian bumps

\textbf{Reproduction plan} (feasibility 5/10, \textasciitilde12h, compute: 1 GPU for RL training + CPU for DG):
\begin{itemize}
\item 1. Build advection HDG solver in MFEM with refinement API.
\item 2. Wrap as Gym environment: obs=element residual features, action=refine/coarsen/nothing.
\item 3. Train DQN per paper Section 4.
\item 4. Compare agent vs bulk(theta=0.5) heuristic on Pareto.
\end{itemize}
\textbf{Expected fidelity:} Qualitative reproduction of Pareto dominance; quantitative within 30\%.. \textbf{Known risks:} Dead code repo → full reimplementation.; HDG implementation details non-trivial..

\textbf{Follow-on questions.}
\begin{itemize}
\item Can the agent generalize across solver types (FVM, DG, CG)?
\item Multi-agent RL for parallel AMR load-balancing.
\item Use of transformer-based policy for long-range correlations.
\item RL-AMR for time-dependent problems with shocks.
\item Offline RL from existing AMR simulation data.
\end{itemize}

\textbf{Rationale for priority:} \emph{Dead code repo. Reimplementation cost is high for what it delivers.}

\medskip\hrule\medskip

\subsection*{\#15. Godunov loss functions for modelling of hyperbolic conservation laws}
\textbf{Authors}: R. Cassia;R. Kerswell \textbullet{} \textbf{Year}: 2024 \textbullet{} \textbf{Cites}: 5 \textbullet{} \textbf{DOI}: \texttt{10.1016/j.cma.2025.117782} \\
\textbf{Priority}: B \textbullet{} \textbf{Revised score}: 13/20 \textbullet{} \textbf{Topic}: Hyperbolic

\textbf{Problem.} Standard PINNs fail for hyperbolic conservation laws with shocks because PDE residual is ill-defined at discontinuities.

\textbf{Approach.} Replace PDE-residual loss with finite-volume-based Godunov/approximate-Riemann-solver flux loss — 'Godunov loss'. Evaluate losses as cell-averaged flux imbalances.

\textbf{Key results.}
\begin{itemize}
\item Beats standard PINN and FVM-based (no-Riemann) losses on 2D Riemann problem.
\item Successful time-stepping and super-resolution reconstruction.
\item Tested across 6 classes of Riemann problems (19 admissible initial conditions).
\end{itemize}

\textbf{Main contribution.} Entropy-consistent shock-capturing loss function for neural PDE solvers.

\textbf{Reproducibility.} Code available: \emph{no}. URL: \texttt{Check CMAME supplementary; search R. Cassia or Kerswell group pages}. Language: python (PyTorch). Completeness: 4/10. Data: 2D Riemann initial conditions (analytical).

\textbf{Benchmarks.} 2D Riemann problem (6 configuration classes)

\textbf{Reproduction plan} (feasibility 6/10, \textasciitilde8h, compute: 1 GPU):
\begin{itemize}
\item 1. Reimplement Godunov loss with HLLC or Roe flux.
\item 2. Generate 2D Riemann initial conditions for one representative configuration per class.
\item 3. Train PINN with Godunov loss; compare to standard PDE loss.
\item 4. Verify sharpness of captured shocks vs TV.
\end{itemize}
\textbf{Expected fidelity:} Qualitative reproduction of shock sharpness; quantitative MSE within 2×.. \textbf{Known risks:} Riemann solver implementation details.; No public code..

\textbf{Follow-on questions.}
\begin{itemize}
\item Extension to 3D Euler?
\item Entropy regularization vs Godunov — ablation.
\item Does WENO loss do better than Godunov loss?
\item Coupling with residual-based adaptive sampling.
\item Performance on ideal MHD?
\end{itemize}

\textbf{Rationale for priority:} \emph{Important topic, tractable reimplementation, but no public code. B tier.}

\medskip\hrule\medskip

\subsection*{\#16. Plane Wave Discontinuous Galerkin Methods for the 2D Helmholtz Equation: Analysis of the p-Version}
\textbf{Authors}: R. Hiptmair;A. Moiola;I. Perugia \textbullet{} \textbf{Year}: 2011 \textbullet{} \textbf{Cites}: 199 \textbullet{} \textbf{DOI}: \texttt{10.1137/090761057} \\
\textbf{Priority}: C \textbullet{} \textbf{Revised score}: 11/20 \textbullet{} \textbf{Topic}: Helmholtz/Wave

\textbf{Problem.} Plane-wave DG methods (Trefftz type) for Helmholtz equation lacked sharp p-version convergence theory.

\textbf{Approach.} Prove p-version convergence estimates for plane-wave DG on 2D Helmholtz; builds on ultra-weak variational formulation.

\textbf{Key results.}
\begin{itemize}
\item Algebraic convergence rate in p with explicit constants.
\item Connection to standard p-FEM theory.
\item Numerical experiments confirming theoretical rates.
\end{itemize}

\textbf{Main contribution.} Rigorous p-version analysis of plane-wave DG for Helmholtz.

\textbf{Reproducibility.} Code available: \emph{no}. URL: \texttt{Not advertised; Hiptmair group ETH may have MATLAB code.}. Language: MATLAB likely. Completeness: 3/10. Data: Analytical Helmholtz solutions on polygonal domains.

\textbf{Benchmarks.} 2D Helmholtz on square with Robin BCs

\textbf{Reproduction plan} (feasibility 5/10, \textasciitilde10h, compute: CPU):
\begin{itemize}
\item 1. Implement plane-wave DG basis in NGSolve.
\item 2. Solve 2D Helmholtz on square for k=1,5,10; vary p.
\item 3. Measure L2 error vs analytical solution; plot convergence in p.
\item 4. Compare observed rate to theory.
\end{itemize}
\textbf{Expected fidelity:} Quantitative match on p-convergence rate.. \textbf{Known risks:} Plane-wave basis has conditioning issues for large p.; No public code..

\textbf{Follow-on questions.}
\begin{itemize}
\item Extension to 3D Helmholtz?
\item Preconditioning of plane-wave DG systems.
\item Coupling with high-order BEM.
\item Time-harmonic elasticity generalization.
\item High-frequency asymptotics connection.
\end{itemize}

\textbf{Rationale for priority:} \emph{Theory-heavy, no code, specialized niche. Demote.}

\medskip\hrule\medskip

\subsection*{\#17. Galerkin Approximations for the Stochastic Burgers Equation}
\textbf{Authors}: D. Blömker;Arnulf Jentzen \textbullet{} \textbf{Year}: 2013 \textbullet{} \textbf{Cites}: 105 \textbullet{} \textbf{DOI}: \texttt{10.1137/110845756} \\
\textbf{Priority}: C \textbullet{} \textbf{Revised score}: 10/20 \textbullet{} \textbf{Topic}: NonlinearScalar

\textbf{Problem.} Rigorous L∞ convergence rates for Galerkin approximations of stochastic Burgers \& related SPDEs.

\textbf{Approach.} Prove spectral Galerkin approximations converge in L∞ (not just L2) for semilinear stochastic evolution equations with additive noise. Apply to stochastic heat, reaction-diffusion, Burgers.

\textbf{Key results.}
\begin{itemize}
\item Explicit rate N\^{}\{-γ\} in truncation with γ tied to noise regularity.
\item L∞ error uniformly in space (new — previously only L2 known).
\item Numerical experiments illustrating rates.
\end{itemize}

\textbf{Main contribution.} First L∞ error estimates for Galerkin-SPDE approximations.

\textbf{Reproducibility.} Code available: \emph{no}. URL: \texttt{No code advertised}. Language: MATLAB likely. Completeness: 2/10. Data: Manufactured SPDE solutions.

\textbf{Benchmarks.} Stochastic heat, Stochastic reaction-diffusion, Stochastic Burgers

\textbf{Reproduction plan} (feasibility 5/10, \textasciitilde8h, compute: CPU):
\begin{itemize}
\item 1. Implement Fourier Galerkin solver for stochastic Burgers.
\item 2. Use Gaussian white-noise Q-Wiener forcing.
\item 3. Monte Carlo over 100+ paths; measure L∞ error.
\item 4. Verify empirical rate matches theoretical γ.
\end{itemize}
\textbf{Expected fidelity:} Quantitative rate match.. \textbf{Known risks:} Monte Carlo variance; need many paths..

\textbf{Follow-on questions.}
\begin{itemize}
\item Sharpness of constants?
\item Extension to multiplicative noise.
\item L∞ rates for finite-element vs spectral.
\item Use in uncertainty quantification.
\item Stochastic Burgers with shocks.
\end{itemize}

\textbf{Rationale for priority:} \emph{Pure-math paper; easy to reimplement numerics but limited 'reproduction' since no code. C tier.}

\medskip\hrule\medskip

\subsection*{\#18. FLUPS: A Fourier-Based Library of Unbounded Poisson Solvers}
\textbf{Authors}: Denis-Gabriel Caprace;Thomas Gillis;P. Chatelain \textbullet{} \textbf{Year}: 2020 \textbullet{} \textbf{Cites}: 18 \textbullet{} \textbf{DOI}: \texttt{10.1137/19M1303848} \\
\textbf{Priority}: B \textbullet{} \textbf{Revised score}: 14/20 \textbullet{} \textbf{Topic}: Elliptic

\textbf{Problem.} Fast Poisson solvers needed for simulations with mixed boundary types (periodic, Neumann, unbounded), particularly vortex methods.

\textbf{Approach.} FLUPS: MPI+OpenMP C++ library, FFT-based Green's function convolution, spectral to order 10; supports all 3x3x3=27 BC combinations. Optimized all-to-all and point-to-point communications.

\textbf{Key results.}
\begin{itemize}
\item Convergence orders 2 to 10+ achieved.
\item Strong scaling to 32K cores on unbounded Poisson.
\item Faster than PMG/FMM at same accuracy for cell-centered unbounded problems.
\end{itemize}

\textbf{Main contribution.} Production-grade MPI Poisson library handling the full BC matrix.

\textbf{Reproducibility.} Code available: \emph{yes}. URL: \texttt{https://gitlab.uliege.be/uliege\_mmech/flups (HTTP 200)}. Language: C++ (MPI + OpenMP, FFTW). Completeness: 9/10. Data: Analytical Poisson test problems in repo.

\textbf{Benchmarks.} Poisson 3D with periodic/Neumann/unbounded BCs, vortex ring

\textbf{Reproduction plan} (feasibility 7/10, \textasciitilde8h, compute: multi-CPU node; HPC for scaling test):
\begin{itemize}
\item 1. Clone flups; build with FFTW.
\item 2. Run convergence test on 3D Poisson; verify order 2,4,6,8,10.
\item 3. Strong-scaling test on 128-1024 cores.
\item 4. Compare timing with PETSc's FFT Poisson.
\end{itemize}
\textbf{Expected fidelity:} Exact match on convergence orders; qualitative match on scaling.. \textbf{Known risks:} FFTW MPI setup.; Access to HPC queue for scaling study..

\textbf{Follow-on questions.}
\begin{itemize}
\item GPU port using cuFFT / heFFTe?
\item Integration with a vortex particle method.
\item Extension to non-uniform grids?
\item Helmholtz variant.
\item Benchmark vs AccFFT.
\end{itemize}

\textbf{Rationale for priority:} \emph{Strong code but HPC-oriented; needs multi-node to really show off. B tier.}

\medskip\hrule\medskip

\subsection*{\#19. Kernel‐based active subspaces with application to computational fluid dynamics parametric problems using the discontinuous Galerkin method}
\textbf{Authors}: F. Romor;M. Tezzele;Andrea Lario;G. Rozza \textbullet{} \textbf{Year}: 2020 \textbullet{} \textbf{Cites}: 15 \textbullet{} \textbf{DOI}: \texttt{10.1002/nme.7099} \\
\textbf{Priority}: B \textbullet{} \textbf{Revised score}: 13/20 \textbullet{} \textbf{Topic}: Hyperbolic

\textbf{Problem.} Active subspaces reduce parameter space linearly; most engineering problems have nonlinear active structures.

\textbf{Approach.} Kernel-based Active Subspaces (KAS) generalize AS with feature maps (Random Fourier Features, radial kernels) and multivariate output objectives. Pipeline with DG-FEM CFD.

\textbf{Key results.}
\begin{itemize}
\item KAS outperforms linear AS on test functions with nonlinear ridge structure.
\item Applied to DG-FEM CFD parametric problem (airfoil).
\item Response-surface accuracy better than Gaussian process baselines.
\end{itemize}

\textbf{Main contribution.} Kernel-AS method for nonlinear dimension reduction in PDE-parametric problems.

\textbf{Reproducibility.} Code available: \emph{yes}. URL: \texttt{https://github.com/mathLab/EZyRB (HTTP 200) + ATHENA from same group}. Language: python. Completeness: 7/10. Data: Benchmark test functions + HopeFOAM CFD runs.

\textbf{Benchmarks.} Radial test functions, SEIR model, DG-FEM airfoil

\textbf{Reproduction plan} (feasibility 7/10, \textasciitilde6h, compute: CPU):
\begin{itemize}
\item 1. Install ATHENA from mathLab.
\item 2. Run radial test function experiments; reproduce linear-vs-kernel AS comparison.
\item 3. (Optional) Run HopeFOAM airfoil case for CFD pipeline.
\end{itemize}
\textbf{Expected fidelity:} Quantitative R² match on test functions.. \textbf{Known risks:} HopeFOAM case requires DG solver stack.; Random seed affects RFF variance..

\textbf{Follow-on questions.}
\begin{itemize}
\item Deep-kernel AS with learned features?
\item Use in Bayesian optimization for expensive PDE surrogates.
\item Extension to multi-fidelity setups.
\item Automatic kernel selection.
\item Active learning of subspace directions.
\end{itemize}

\textbf{Rationale for priority:} \emph{Tool-chain heavier than advertised; fine B-tier candidate for UQ/surrogates.}

\medskip\hrule\medskip

\subsection*{\#20. On the convergence of discontinuous Galerkin/Hermite spectral methods for the Vlasov-Poisson system}
\textbf{Authors}: M. Bessemoulin-Chatard;F. Filbet \textbullet{} \textbf{Year}: 2022 \textbullet{} \textbf{Cites}: 6 \textbullet{} \textbf{DOI}: \texttt{10.48550/arXiv.2208.12503} \\
\textbf{Priority}: C \textbullet{} \textbf{Revised score}: 10/20 \textbullet{} \textbf{Topic}: Elliptic

\textbf{Problem.} DG/Hermite-spectral methods for Vlasov-Poisson lacked a priori convergence theory.

\textbf{Approach.} Use asymmetrically weighted Hermite in velocity + DG in space. Introduce time-dependent weighted L2 space. Prove stability + error estimate.

\textbf{Key results.}
\begin{itemize}
\item Global weighted L2 stability.
\item Propagation of regularity.
\item Error estimate in weighted L2.
\end{itemize}

\textbf{Main contribution.} First convergence theory for DG/Hermite Vlasov-Poisson.

\textbf{Reproducibility.} Code available: \emph{no}. URL: \texttt{No code advertised}. Language: Julia/Fortran likely. Completeness: 2/10. Data: Manufactured Vlasov-Poisson solutions.

\textbf{Benchmarks.} Landau damping, two-stream instability, bump-on-tail

\textbf{Reproduction plan} (feasibility 4/10, \textasciitilde12h, compute: CPU):
\begin{itemize}
\item 1. Implement DG-Hermite Vlasov-Poisson from scratch in Python + FEniCSx.
\item 2. Landau damping test; verify damping rate matches linear theory.
\item 3. Refine Hermite moments; measure convergence.
\item 4. Compare to reported rate.
\end{itemize}
\textbf{Expected fidelity:} Qualitative match; quantitative rate within 20\%.. \textbf{Known risks:} Significant implementation effort; weighted-Hermite basis conditioning..

\textbf{Follow-on questions.}
\begin{itemize}
\item Extension to Vlasov-Maxwell?
\item Positivity-preserving DG-Hermite schemes.
\item Kinetic plasma turbulence benchmarks.
\item Comparison to semi-Lagrangian.
\item GPU port.
\end{itemize}

\textbf{Rationale for priority:} \emph{Theory paper; heavy reimplementation, no code. C tier.}

\medskip\hrule\medskip

\subsection*{\#21. Improvements to the APBS biomolecular solvation software suite}
\textbf{Authors}: E. Jurrus;David Engel;Keith T. Star;Kyle Monson;Juan Brandi;Lisa E. Felberg;David H. Brookes;Leighton Wilson;Jiahui Chen;Karina Liles;Minju Chun;Peter Li;D. Gohara;Todd J. Dolinsky;R. Konečný;D. Koes; \textbullet{} \textbf{Year}: 2017 \textbullet{} \textbf{Cites}: 2081 \textbullet{} \textbf{DOI}: \texttt{10.1002/pro.3280} \\
\textbf{Priority}: A \textbullet{} \textbf{Revised score}: 15/20 \textbullet{} \textbf{Topic}: Elliptic

\textbf{Problem.} Update and extend APBS for biomolecular electrostatics (Poisson-Boltzmann equation).

\textbf{Approach.} New features in APBS suite: analytical/semi-analytical PBE solver, boundary element solver, geometric flow solvation, graph-based pKa, visualization.

\textbf{Key results.}
\begin{itemize}
\item Scalable to large biomolecules.
\item Multiple solver backends (FD multigrid, BEM, analytical).
\item Validated on standard benchmark sets.
\end{itemize}

\textbf{Main contribution.} Reference biomolecular electrostatics toolkit, widely used (2081 cites).

\textbf{Reproducibility.} Code available: \emph{yes}. URL: \texttt{https://github.com/Electrostatics/apbs (HTTP 200)}. Language: C/C++ with Python bindings. Completeness: 10/10. Data: PQR files from PDB, provided in docs.

\textbf{Benchmarks.} alpha-chymotrypsin, insulin hexamer, lysozyme, 10+ standard biomol cases

\textbf{Reproduction plan} (feasibility 9/10, \textasciitilde6h, compute: CPU (multi-core)):
\begin{itemize}
\item 1. conda install -c conda-forge apbs pdb2pqr.
\item 2. Run example inputs from APBS tutorial (1ake lysozyme etc).
\item 3. Reproduce solvation energy values.
\item 4. Use geometric-flow solver on one test; compare to FD-multigrid.
\end{itemize}
\textbf{Expected fidelity:} Quantitative match to within 1-2\%.. \textbf{Known risks:} Compiling APBS from source is painful; conda package is stable..

\textbf{Follow-on questions.}
\begin{itemize}
\item Integration with ML potentials (FEP/ML).
\item Coupling with MD via AMBER PBSA.
\item GPU port of geometric flow.
\item AlphaFold → APBS automated pipeline.
\item pKa prediction benchmark vs PROPKA.
\end{itemize}

\textbf{Rationale for priority:} \emph{Mature code, easy conda install. Upgrade — a very safe starter replication.}

\medskip\hrule\medskip

\subsection*{\#22. Fast Nonsymmetric Iterations and Preconditioning for Navier-Stokes Equations}
\textbf{Authors}: H. Elman;D. Silvester \textbullet{} \textbf{Year}: 1996 \textbullet{} \textbf{Cites}: 258 \textbullet{} \textbf{DOI}: \texttt{10.1137/0917004} \\
\textbf{Priority}: B \textbullet{} \textbf{Revised score}: 14/20 \textbullet{} \textbf{Topic}: Fluid/NavierStokes

\textbf{Problem.} Classic Navier-Stokes linear systems are non-symmetric + indefinite; few preconditioners give mesh-independent convergence.

\textbf{Approach.} Block preconditioners for discrete Oseen operator: diagonal with Schur-complement approximations. Prove eigenvalues bounded independent of h.

\textbf{Key results.}
\begin{itemize}
\item GMRES iterations independent of mesh size h.
\item Inner iterative solvers for preconditioner keep costs low.
\item Validated on backward-facing step and lid-driven cavity.
\end{itemize}

\textbf{Main contribution.} Foundational paper (258 cites) for block preconditioning of incompressible Navier-Stokes — precursor to PCD/LSC.

\textbf{Reproducibility.} Code available: \emph{no}. URL: \texttt{None advertised (1996 paper); modern equivalents in IFISS (MATLAB)}. Language: MATLAB. Completeness: 6/10. Data: Standard Stokes/Oseen benchmarks.

\textbf{Benchmarks.} Backward-facing step, lid-driven cavity, channel flow

\textbf{Reproduction plan} (feasibility 8/10, \textasciitilde5h, compute: CPU):
\begin{itemize}
\item 1. Download IFISS; run navier\_testproblem with the block preconditioner.
\item 2. Reproduce iteration-vs-h table at multiple Re.
\item 3. Verify mesh-independence claim.
\end{itemize}
\textbf{Expected fidelity:} Quantitative iteration-count reproduction.. \textbf{Known risks:} MATLAB license for IFISS.; Original 1996 numbers may differ slightly from IFISS defaults..

\textbf{Follow-on questions.}
\begin{itemize}
\item Compare to LSC and PCD preconditioners.
\item Extend to stabilized FE discretizations.
\item Use in ROM surrogates.
\item GPU port of block preconditioner.
\item Extension to Navier-Stokes-Boussinesq.
\end{itemize}

\textbf{Rationale for priority:} \emph{Classic, reproducible via IFISS. Solid B.}

\medskip\hrule\medskip

\subsection*{\#23. New Laplace and Helmholtz solvers}
\textbf{Authors}: A. Gopal;L. Trefethen \textbullet{} \textbf{Year}: 2019 \textbullet{} \textbf{Cites}: 51 \textbullet{} \textbf{DOI}: \texttt{10.1073/pnas.1904139116} \\
\textbf{Priority}: A \textbullet{} \textbf{Revised score}: 15/20 \textbullet{} \textbf{Topic}: Helmholtz/Wave

\textbf{Problem.} Standard FEM/BEM methods have O(h\^{}p) polynomial convergence on polygons with corners.

\textbf{Approach.} Rational-function approximations (lightning solvers) with exponential clustering of poles near corners give fractional-exponential (root-exponential) convergence for 2D Laplace/Helmholtz.

\textbf{Key results.}
\begin{itemize}
\item e\^{}\{-c√N\} convergence for 2D Laplace on polygons with corners.
\item Orders of magnitude faster than FEM/BEM for comparable accuracy.
\item Implemented in lightning.m (MATLAB).
\end{itemize}

\textbf{Main contribution.} Paradigm shift toward rational-function PDE solvers (lightning Laplace solver).

\textbf{Reproducibility.} Code available: \emph{yes}. URL: \texttt{https://people.maths.ox.ac.uk/trefethen/lightning.html (MATLAB scripts)}. Language: MATLAB. Completeness: 9/10. Data: Analytical Laplace problems on L-shape, polygons.

\textbf{Benchmarks.} Laplace on L-shape, Helmholtz on polygon with corners

\textbf{Reproduction plan} (feasibility 9/10, \textasciitilde3h, compute: CPU (laptop)):
\begin{itemize}
\item 1. Download lightning.m from Trefethen's page.
\item 2. Run L-shape example; verify 10-digit accuracy.
\item 3. Vary N; plot convergence vs e\^{}\{-c√N\}.
\item 4. Compare time-to-10-digits vs FEniCS on same geometry.
\end{itemize}
\textbf{Expected fidelity:} Exact match at machine precision.. \textbf{Known risks:} MATLAB license..

\textbf{Follow-on questions.}
\begin{itemize}
\item Python port (already exists via Baddoo's lightning-python)?
\item 3D extension?
\item Lightning Helmholtz at high k.
\item Combining lightning with FEM (hybrid)?
\item Use in conformal mapping.
\end{itemize}

\textbf{Rationale for priority:} \emph{Short, clean, high-impact paper with a one-file MATLAB demo. Perfect reproduction target.}

\medskip\hrule\medskip

\subsection*{\#24. Stability and error analysis of a class of high-order IMEX schemes for Navier-stokes equations with periodic boundary conditions}
\textbf{Authors}: Fukeng Huang;Jie Shen \textbullet{} \textbf{Year}: 2021 \textbullet{} \textbf{Cites}: 43 \textbullet{} \textbf{DOI}: \texttt{10.1137/21m1404144} \\
\textbf{Priority}: B \textbullet{} \textbf{Revised score}: 12/20 \textbullet{} \textbf{Topic}: Fluid/NavierStokes

\textbf{Problem.} High-order IMEX schemes for Navier-Stokes lacked unconditional-stability error analysis.

\textbf{Approach.} Scalar-auxiliary-variable (SAV) technique + IMEX k-step (k=1..5) + Fourier-Galerkin space. Prove unconditional L∞-H\^{}1 stability and error estimates.

\textbf{Key results.}
\begin{itemize}
\item First unconditionally stable high-order IMEX for Navier-Stokes.
\item Up to 5th-order in time, confirmed numerically.
\item Successful on double-shear-layer problem.
\end{itemize}

\textbf{Main contribution.} Rigorous high-order IMEX-SAV framework for incompressible NS.

\textbf{Reproducibility.} Code available: \emph{no}. URL: \texttt{Shen group sometimes publishes MATLAB on personal pages — check}. Language: MATLAB likely. Completeness: 3/10. Data: Manufactured solutions + double shear layer IC.

\textbf{Benchmarks.} Taylor-Green 2D, double shear layer 2D

\textbf{Reproduction plan} (feasibility 6/10, \textasciitilde8h, compute: CPU):
\begin{itemize}
\item 1. Implement Fourier pseudo-spectral Navier-Stokes in 2D.
\item 2. Add SAV auxiliary variable; implement IMEX-BDFk for k=1..5.
\item 3. Manufactured-solution test; verify order k in Δt.
\item 4. Double shear layer at Re=10000; compare vorticity to reported figure.
\end{itemize}
\textbf{Expected fidelity:} Quantitative order match; qualitative vorticity match.. \textbf{Known risks:} SAV implementation subtleties.; High-order IMEX conditioning..

\textbf{Follow-on questions.}
\begin{itemize}
\item SAV + exponential integrators?
\item Extension to MHD.
\item SAV on unstructured FE meshes.
\item Combination with adaptive time-stepping.
\item Parallelization.
\end{itemize}

\textbf{Rationale for priority:} \emph{Clear numerics but no code; reimplementation tractable.}

\medskip\hrule\medskip

\subsection*{\#25. Rate optimality of adaptive finite element methods with respect to the overall computational costs}
\textbf{Authors}: G. Gantner;A. Haberl;D. Praetorius;S. Schimanko \textbullet{} \textbf{Year}: 2020 \textbullet{} \textbf{Cites}: 34 \textbullet{} \textbf{DOI}: \texttt{10.1090/mcom/3654} \\
\textbf{Priority}: C \textbullet{} \textbf{Revised score}: 11/20 \textbullet{} \textbf{Topic}: Elliptic

\textbf{Problem.} Adaptive FEM rate-optimality theory typically measured vs \# DOF, not vs computational cost.

\textbf{Approach.} Analyze adaptive FEM with inexact iterative solvers (CG for linear, Zarantonello for nonlinear monotone). Prove quasi-optimal linear convergence w.r.t. total cost.

\textbf{Key results.}
\begin{itemize}
\item Rate-optimality in overall cost (not just DOF).
\item Applies to linear elliptic and nonlinear monotone problems.
\item Numerical experiments confirm.
\end{itemize}

\textbf{Main contribution.} First rate-optimality theorem w.r.t. overall computational cost for inexact AFEM.

\textbf{Reproducibility.} Code available: \emph{no}. URL: \texttt{TU Wien group sometimes publishes MATLAB}. Language: MATLAB likely. Completeness: 3/10. Data: Standard elliptic test problems.

\textbf{Benchmarks.} L-shape Laplace, Z-shape elliptic, p-Laplacian nonlinear

\textbf{Reproduction plan} (feasibility 6/10, \textasciitilde8h, compute: CPU):
\begin{itemize}
\item 1. Implement residual-based AFEM in FEniCS with Dörfler marking.
\item 2. Use inexact CG for linear system; adjust solver tolerance adaptively.
\item 3. Track error vs cumulative solver cost on L-shape.
\item 4. Confirm algebraic rate -1 w.r.t. cost.
\end{itemize}
\textbf{Expected fidelity:} Qualitative reproduction of rate.. \textbf{Known risks:} AFEM implementation details subtle.; No public reference code..

\textbf{Follow-on questions.}
\begin{itemize}
\item Extension to parabolic problems.
\item Goal-oriented AFEM cost-optimality.
\item Practical comparison w/ heuristic AFEM.
\item Multigrid inside inexact solver.
\item ML-based marking strategies.
\end{itemize}

\textbf{Rationale for priority:} \emph{Theory-focused, no code. C.}

\medskip\hrule\medskip

\subsection*{\#26. Comparison of Adaptive Multiresolution and Adaptive Mesh Refinement Applied to Simulations of the Compressible Euler Equations}
\textbf{Authors}: R. Deiterding;M. O. Domingues;S. Gomes;K. Schneider \textbullet{} \textbf{Year}: 2015 \textbullet{} \textbf{Cites}: 33 \textbullet{} \textbf{DOI}: \texttt{10.1137/15M1026043} \\
\textbf{Priority}: B \textbullet{} \textbf{Revised score}: 13/20 \textbullet{} \textbf{Topic}: AMR/Multigrid

\textbf{Problem.} AMR and adaptive-multiresolution (MR) are two dominant approaches for PDEs with sharp fronts; which is better?

\textbf{Approach.} Head-to-head FV comparison on compressible Euler (2D Riemann Lax-Liu \#6, 3D expanding shock). Both use 2nd-order shock-capturing + local time-stepping.

\textbf{Key results.}
\begin{itemize}
\item MR has better memory compression and slightly enhanced convergence.
\item AMR has lower absolute overhead in tested codes.
\item CPU-time compression trends similar for both.
\end{itemize}

\textbf{Main contribution.} Quantitative comparison settling MR vs AMR tradeoffs on standard Euler benchmarks.

\textbf{Reproducibility.} Code available: \emph{yes}. URL: \texttt{https://github.com/waveletApplications/carmen (HTTP 200) for MR side; AMROC for AMR}. Language: C/C++/Fortran. Completeness: 7/10. Data: Riemann IC analytical.

\textbf{Benchmarks.} 2D Riemann Lax-Liu \#6, 3D expanding shock ellipsoid

\textbf{Reproduction plan} (feasibility 7/10, \textasciitilde8h, compute: multi-CPU):
\begin{itemize}
\item 1. Build Carmen; run Lax-Liu \#6.
\item 2. Build AMROC; run same problem.
\item 3. Compare memory and CPU vs regular-grid baseline.
\item 4. Confirm MR memory advantage.
\end{itemize}
\textbf{Expected fidelity:} Quantitative compression factor match within 20\%.. \textbf{Known risks:} AMROC build is specialized.; Timings hardware-dependent..

\textbf{Follow-on questions.}
\begin{itemize}
\item Hybrid MR+AMR?
\item Performance on 3D turbulence.
\item GPU AMR frameworks (AMReX) comparison.
\item Wavelet-based refinement indicator portability.
\item Extension to relativistic hydro.
\end{itemize}

\textbf{Rationale for priority:} \emph{Both codes exist but heavy-weight; mid-tier B.}

\medskip\hrule\medskip

\subsection*{\#27. Kinetic.jl: A portable finite volume toolbox for scientific and neural computing}
\textbf{Authors}: Tianbai Xiao \textbullet{} \textbf{Year}: 2021 \textbullet{} \textbf{Cites}: 17 \textbullet{} \textbf{DOI}: \texttt{10.21105/joss.03060} \\
\textbf{Priority}: A \textbullet{} \textbf{Revised score}: 15/20 \textbullet{} \textbf{Topic}: ML-operator/PINN

\textbf{Problem.} Need a flexible FV framework for kinetic equations + scientific ML in Julia.

\textbf{Approach.} Kinetic.jl: modular FV toolbox with KitBase (physics) and KitML (neural closures). Supports multi-scale Boltzmann-Navier-Stokes transitions.

\textbf{Key results.}
\begin{itemize}
\item Solves Navier-Stokes and Boltzmann in unified API.
\item Neural closure example for BGK collision operator.
\item Short JOSS paper (software publication).
\end{itemize}

\textbf{Main contribution.} Unified kinetic-theory + ML Julia toolbox.

\textbf{Reproducibility.} Code available: \emph{yes}. URL: \texttt{https://github.com/vavrines/Kinetic.jl (HTTP 200)}. Language: Julia. Completeness: 9/10. Data: Test scripts in repo.

\textbf{Benchmarks.} Sod shock tube, Couette flow, BGK Boltzmann

\textbf{Reproduction plan} (feasibility 9/10, \textasciitilde4h, compute: CPU):
\begin{itemize}
\item 1. Install Julia 1.10+; add Kinetic.jl.
\item 2. Run examples/sod.jl; verify shock tube matches analytical Riemann solution.
\item 3. Run Couette flow.
\item 4. Reproduce a neural-closure example from KitML.
\end{itemize}
\textbf{Expected fidelity:} Quantitative match on shock positions/states.. \textbf{Known risks:} Julia version compatibility with older package commits..

\textbf{Follow-on questions.}
\begin{itemize}
\item Multi-scale coupling between kinetic and NS regimes.
\item GPU support via CUDA.jl.
\item Plasma kinetics.
\item Neural BGK operator learning benchmarks.
\item Adaptive grid refinement inside Kinetic.jl.
\end{itemize}

\textbf{Rationale for priority:} \emph{Julia ease + short paper = very low friction A-tier replication.}

\medskip\hrule\medskip

\subsection*{\#28. Highly efficient and energy dissipative schemes for the time fractional Allen-Cahn equation}
\textbf{Authors}: Dianming Hou;Chuanju Xu \textbullet{} \textbf{Year}: 2021 \textbullet{} \textbf{Cites}: 17 \textbullet{} \textbf{DOI}: \texttt{10.1137/20m135577x} \\
\textbf{Priority}: C \textbullet{} \textbf{Revised score}: 10/20 \textbullet{} \textbf{Topic}: PhaseField/RD

\textbf{Problem.} Time-fractional Allen-Cahn equation needs stable schemes on graded meshes for efficient long-time simulation.

\textbf{Approach.} Split fractional derivative into local + history. Use L1/L1-CN/L1+-CN for local; extended SAV for nonlinearity + history. Fourier-Galerkin space.

\textbf{Key results.}
\begin{itemize}
\item Unconditional energy dissipation proven.
\item Dissipation holds on graded mesh.
\item Numerical confirmation on 2D Allen-Cahn.
\end{itemize}

\textbf{Main contribution.} First unconditionally dissipative high-order scheme for time-fractional Allen-Cahn.

\textbf{Reproducibility.} Code available: \emph{no}. URL: \texttt{No code advertised}. Language: MATLAB likely. Completeness: 2/10. Data: Manufactured + standard Allen-Cahn IC.

\textbf{Benchmarks.} Time-fractional Allen-Cahn 2D

\textbf{Reproduction plan} (feasibility 5/10, \textasciitilde10h, compute: CPU):
\begin{itemize}
\item 1. Implement L1+-CN on graded mesh with sum-of-exponentials history.
\item 2. Manufactured solution test for convergence order.
\item 3. 2D Allen-Cahn coarsening; verify energy monotonicity.
\end{itemize}
\textbf{Expected fidelity:} Quantitative order match.. \textbf{Known risks:} Sum-of-exponentials approximation accuracy.; No reference code..

\textbf{Follow-on questions.}
\begin{itemize}
\item Extension to time-fractional Cahn-Hilliard.
\item Higher-order L2/L3 schemes.
\item Parallel SoE history evaluation.
\item Coupling to phase-field crystal models.
\item Adaptive time-stepping.
\end{itemize}

\textbf{Rationale for priority:} \emph{No code; specialized.}

\medskip\hrule\medskip

\subsection*{\#29. Modified Poisson-Nernst-Planck Model with Coulomb and Hard-sphere Correlations}
\textbf{Authors}: Manman Ma;Zhenli Xu;Liwei Zhang \textbullet{} \textbf{Year}: 2020 \textbullet{} \textbf{Cites}: 15 \textbullet{} \textbf{DOI}: \texttt{10.1137/19m1310098} \\
\textbf{Priority}: C \textbullet{} \textbf{Revised score}: 10/20 \textbullet{} \textbf{Topic}: Elliptic

\textbf{Problem.} Standard Poisson-Nernst-Planck misses Coulomb and hard-sphere correlations important for multivalent ions at interfaces.

\textbf{Approach.} Add Coulomb correlation via generalized Debye-Hückel (self-Green's function) and hard-sphere correlation via modified fundamental measure theory. Solve via WKB asymptotics for parallel-plate geometry.

\textbf{Key results.}
\begin{itemize}
\item Reproduces charge-inversion effects seen in Monte Carlo simulations.
\item Agreement with MC for multivalent ions.
\item Analytical WKB solution for two-plate case.
\end{itemize}

\textbf{Main contribution.} Modified PNP with both long- and short-range corrections beyond mean-field.

\textbf{Reproducibility.} Code available: \emph{no}. URL: \texttt{No code advertised}. Language: MATLAB likely. Completeness: 2/10. Data: Reference MC data from cited papers.

\textbf{Benchmarks.} Parallel-plate electrolyte with multivalent ions

\textbf{Reproduction plan} (feasibility 5/10, \textasciitilde10h, compute: CPU):
\begin{itemize}
\item 1. Implement generalized Debye-Hückel equation solver (1D).
\item 2. Add MFMT hard-sphere correction.
\item 3. WKB asymptotic solution for two plates.
\item 4. Compare density profiles to reference MC data.
\end{itemize}
\textbf{Expected fidelity:} Qualitative match on charge inversion.. \textbf{Known risks:} Implementation of self-Green's-function iteration.; No reference code..

\textbf{Follow-on questions.}
\begin{itemize}
\item Extension to 3D finite systems.
\item Coupling with ML for closure.
\item Nonequilibrium transport predictions.
\item Nanopore applications.
\item Benchmark vs classical DFT of ions.
\end{itemize}

\textbf{Rationale for priority:} \emph{Niche physics; no code.}

\medskip\hrule\medskip

\subsection*{\#30. Optimized Schwarz Methods without Overlap for the Helmholtz Equation}
\textbf{Authors}: M. Gander;F. Magoulès;F. Nataf \textbullet{} \textbf{Year}: 2002 \textbullet{} \textbf{Cites}: 403 \textbullet{} \textbf{DOI}: \texttt{10.1137/S1064827501387012} \\
\textbf{Priority}: B \textbullet{} \textbf{Revised score}: 13/20 \textbullet{} \textbf{Topic}: Helmholtz/Wave

\textbf{Problem.} Classical Schwarz for Helmholtz fails without overlap (propagating modes undamped).

\textbf{Approach.} Derive optimized Robin-type transmission conditions (symbol σ(k) optimized in Fourier). Proves convergence without overlap for Helmholtz.

\textbf{Key results.}
\begin{itemize}
\item Optimal transmission condition gives convergence in finite steps (nonlocal) + local approximation.
\item Dramatic iteration-count reduction over classical Schwarz.
\item Numerical validation on 2D Helmholtz benchmark.
\end{itemize}

\textbf{Main contribution.} Foundational (403 cites) optimized Schwarz theory for Helmholtz.

\textbf{Reproducibility.} Code available: \emph{no}. URL: \texttt{None advertised (2002 paper); modern HPDDM library includes OSM}. Language: MATLAB / FreeFem++. Completeness: 6/10. Data: Analytical Helmholtz solutions.

\textbf{Benchmarks.} 2D Helmholtz in rectangle

\textbf{Reproduction plan} (feasibility 7/10, \textasciitilde5h, compute: CPU):
\begin{itemize}
\item 1. Implement Schwarz iteration with classical vs optimized Robin BC in FreeFem++.
\item 2. 2D Helmholtz on [0,1]² at k=10.
\item 3. Compare iteration counts.
\item 4. Reproduce convergence-rate curves from paper.
\end{itemize}
\textbf{Expected fidelity:} Quantitative iteration-count reproduction.. \textbf{Known risks:} Tuning of σ parameter.; 2002 numerics environment..

\textbf{Follow-on questions.}
\begin{itemize}
\item Extension to Maxwell's equations.
\item Hybrid OSM + sweeping preconditioner.
\item OSM on unstructured meshes.
\item Performance at very high k.
\item ML-tuning of optimized parameters.
\end{itemize}

\textbf{Rationale for priority:} \emph{Classic paper reproducible via HPDDM or direct FreeFem++ implementation.}

\medskip\hrule\medskip

\section{Recommended 15-paper batch schedule}

See companion \texttt{top30\_batch\_schedule.md} for details. Summary:

\textbf{Wave 1} (5 parallel, \textasciitilde 6--8h wallclock, CPU+GPU mix): FEM-vs-PINNs, Walk on Stars, Fast Poisson Spectral, Lightning Laplace, Kinetic.jl.

\textbf{Wave 2} (5 parallel, \textasciitilde 6--10h, 1 GPU each): jax-cfd (ML-CFD), Latent Spectral Models, Koopman Neural Operator, Laplace Neural Operator, Dedalus.

\textbf{Wave 3} (5 parallel, mostly CPU-heavy): APBS, lifex-cfd, Elman-Silvester (IFISS), Optimized Schwarz, FLUPS.

Total wallclock with 5-way parallelism: \textasciitilde 24--30 hours.

\section{Conclusion}

Rick: our highest-yield near-term set is the 11 Priority~A papers; starting Wave~1 of 5 today produces verifiable results in a workday. Stage-5 demoted 5 papers to C due to dead/absent code; Stage-4 signal heuristics (which scored \texttt{github:...} substring presence without HTTP validation) were too optimistic on those. One 404 (Foucart 2022 RL-AMR) is a clean red flag we couldn't have caught without this stage.

\end{document}
